\documentclass[12pt]{article}

\newif\ifsol
\soltrue

% Unicode compatibility
\usepackage{iftex}
\ifPDFTeX
  \usepackage[utf8]{inputenc}
  \usepackage[noTeX]{mmap}
  \usepackage[T1]{fontenc}
\fi
% XeTeX does not support mmap
\ifLuaTeX
  \usepackage{luatex85}
  \usepackage[noTeX]{mmap}
\fi

\setlength{\textheight}{9in}
\setlength{\textwidth}{7.1in}
\setlength{\evensidemargin}{-0.2in}
\setlength{\oddsidemargin}{-0.2in}
\setlength{\headsep}{30pt}
\setlength{\topmargin}{-0.3in}

\usepackage{amsthm,amsfonts,amsmath,xspace,tikz,varwidth,cancel,soul}
\usepackage{hyperref,verbatim}
\newtheorem{theorem}{Theorem}
\theoremstyle{definition}
\newtheorem{definition}{Definition}

\newcommand{\addmedskip}{\addvspace{\medskipamount}}

\newcommand{\addbigskip}{\addvspace{\bigskipamount}}

\pagestyle{myheadings}
\markboth{BU CAS CS 538.}{BU CAS CS 538.} 
\newcounter{problemnum}
\setcounter{problemnum}{0}
\newenvironment{problem}
     {\addbigskip \setcounter{partnum}{0}
      \noindent\stepcounter{problemnum}\textbf{Problem
                                             \arabic{problemnum}.\ }}
     {\par\addbigskip}
     
\ifsol
\newenvironment{solution}
     {\addbigskip
      \noindent\color{blue}\textbf{Solution.\ }}
     {\par\addbigskip}
\else
\newenvironment{solution}
     {\comment}
     {\endcomment}
\fi

\newcounter{partnum}
\setcounter{partnum}{0}
\newenvironment{ppart}
     {\addmedskip
      \noindent\stepcounter{partnum}\textbf{(\alph{partnum})}\ }
     {\par\addbigskip}

\newcommand{\Enc}{\mathrm{Enc}}
\newcommand{\Dec}{\mathrm{Dec}}
\newcommand{\pk}{\mathit{pk}}
\newcommand{\sk}{\mathit{sk}}
\newcommand{\qrn}{\mathit{QR}_n}
\newcommand{\qrp}{\mathit{QR}_p}
\newcommand{\qrpone}{\mathit{QR}_{p_1}}
\newcommand{\qrptwo}{\mathit{QR}_{p_2}}
\newcommand{\prg}{\mathit{J}}

\newcommand{\crt}{\mathrm{crt}}


\newcommand{\eqdef}{\buildrel{\scriptstyle{\mathit{def}}}\over{=}}
\newcommand{\rfrom}{\buildrel{\scriptstyle{\mathrm{R}}}\over{\leftarrow}}

\newcommand{\mbmod}{\,\%\,}

\newcommand{\A}{\mathcal{A}}
\renewcommand{\L}{\mathcal{L}}
\newcommand{\Z}{\mathbb{Z}}
\newcommand{\G}{\mathbb{G}}
\newcommand{\X}{\mathcal{X}}
\newcommand{\Y}{\mathcal{Y}}
\newcommand{\K}{\mathcal{K}}
\newcommand{\M}{\mathcal{M}}
\newcommand{\E}{\mathcal{E}}
\newcommand{\C}{\mathcal{C}}
\newcommand{\R}{\mathcal{R}}
\renewcommand{\S}{\mathcal{S}}
\newcommand{\I}{\mathcal{I}}
\newcommand{\T}{\mathcal{T}}
\newcommand{\randk}{\mathsf k}
\newcommand{\randm}{\mathsf m}
\newcommand{\randc}{\mathsf c}
\newcommand{\Adv}{\mathcal{A}}
\newcommand{\B}{\mathcal{B}}
\newcommand{\SSadv}{\textrm{SS}\textsf{adv}}
\newcommand{\PRGadv}{\textrm{PRG}\textsf{adv}}
\newcommand{\threebitadv}{\textrm{3Bit}\textsf{adv}}
\newcommand{\PRFadv}{\textrm{PRF}\textsf{adv}}
\newcommand{\MACadv}{\textrm{MAC}\textsf{adv}}
\newcommand{\CPAadv}{\textrm{CPA}\textsf{adv}}
\newcommand{\BCadv}{\textrm{BC}\textsf{adv}}
\newcommand{\CRadv}{\textrm{CR}\textsf{adv}}
\newcommand{\DDHprimeadv}{\textrm{DDH}'\textsf{adv}}
\newcommand{\DDHadv}{\textrm{DDH}\textsf{adv}}
\newcommand{\DDHprime}{DDH$'$}

\newtheorem{lemma}{Lemma}

\newcommand{\indist}{  \mathrel{\vcenter{\offinterlineskip
  \hbox{$\sim$}\vskip-.35ex\hbox{$\sim$}\vskip-.35ex\hbox{$\sim$}}}}

\newcommand{\samples}{\xleftarrow{R}}

\newcommand{\hdl}{H_{\mathrm{dl}}}

%%% BEGIN MIKE ROSULEK'S MACROS
\newcommand{\codebox}[1]{%
        \begin{varwidth}{\linewidth}%
        \begin{tabbing}%
            ~~~\=\quad\=\quad\=\quad\=\kill
            #1
        \end{tabbing}%
        \end{varwidth}%
}

\newcommand{\fcodebox}[1]{%
    \fcolorbox{black}{white}{\codebox{#1}}%
}

\newcommand{\titlecodebox}[2]{%
    \fboxsep=0pt%
    \fcolorbox{black}{black!10}{%
        \begin{varwidth}{\linewidth}%
        \centering%
        \fboxsep=3pt%
        \colorbox{black!10}{#1} \\
        \colorbox{white}{\codebox{#2}}%
        \end{varwidth}%
    }
}



\newcommand{\link}{\diamond}
\newcommand{\lib}[1]{\ensuremath{\L_{\textsf{#1}}}\xspace}
\newcommand{\subname}[1]{\ensuremath{\textsc{#1}}\xspace}

%%% END OF MIKE ROSULEK'S MACROS


\begin{document}
\begin{center}
\Large{\textbf{CAS CS 538.  \ifsol Solutions to \fi Problem Set 8}}\\
\smallskip
\large{\textbf{Due electronically via Gradescope, Monday April 6, 2026 11:59pm}}//
\large{\textbf{Om Khadka, U51801771}}
\end{center}

\noindent
\textbf{How to submit:}  See instructions on Problem Set 1. 
 


 \paragraph{Trapdoor Hash Functions.}


 \noindent
For this problem, we will augment collision-resistant hashing with a useful property.
 
We will consider keyed collision-resistant hash functions that take a hashing key $\pk$ and two input values (which we will call $(a, b)$), and output one value. We will now augment these hash functions with an additional property: a trapdoor secret key $\sk$ that enables easy collision finding. 

We formalize this property as follows. We add a key generation algorithm
 $G$ that outputs a public hashing key $\pk$ and a trapdoor key $\sk$. There is also an efficient algorithm
 $T$ to find collisions given $\sk$: specifically,  $T(\sk, a_1, b_1, a_2)$, outputs $b_2$ such
 that $H(\pk, a_1, b_1)=H(\pk, a_2, b_2)$.  We still require that the hash function remains collision resistant for anyone who is given $\pk$ but not $\sk$.


 In the next two problems you will build trapdoor hash functions. Of course, if a hash function has a trapdoor, then the person who generated the key will be able to find collisions.
 
 \begin{problem} (30 points) \label{dltrapdoor}
Recall the DL-based family in Section 8.5.1 of the textbook. Let $q$ be a Sophie Germain prime, so that $p=2q+1$ is also prime. Let  $x, y\in\qrp$ and $a, b\in \{1, 2, \dots, q\}$. The public key is $(x, y)$ and $H((x, y), a, b)=\mathrm{abs} (x^a y^b\bmod p)$, where $\mathrm{abs}(z)$ returns $z$ if $z<p/2$ and $-z$ otherwise.

Show that $H$ is actually a trapdoor hash
 family.  In other words, demonstrate an algorithm $G$ creates $\pk=(x, y)$ and some $\sk$ (your job is to figure out what $\sk$ should be), and algorithm $T$.
 \end{problem}

\begin{solution}
Your solution goes here
\end{solution}


 \begin{problem}  (40 points at 10 each)
  In this problem, we will build a collision-resistant hash function based on the hardness of factoring rather on the discrete logarithm problem. 
 
 Let $q_1$ and $q_2$ be two distinct Sophie Germain primes, so that $p_1=2q_1+1$ and $p_2=2q_2+1$ are primes. Let $n=p_1 p_2$. Denote by $\qrpone$, $\qrptwo$, and $\qrn$ the groups of squares modulo $p_1$, $p_2$, and $n$, respectively. We know that $\qrpone$ and $\qrptwo$ are cyclic; thus, let $g_1\in \qrpone$ generate $\qrpone$ and $g_2\in\qrpone$ generate $\qrptwo$.

 \begin{ppart}
  Show by CRT (see notes from Discussion 8) that $\qrn$ is cyclic (the easiest way to do this is to try to construct a generator of $\qrn$ and show that it is indeed a generator by showing that its order is $q_1q_2$).  
 \end{ppart}

\begin{solution}
Your solution goes here
\end{solution}
 

 \begin{ppart} 
 Show how to efficiently factor $n$ given $q_1 q_2$. In fact, you can factor $n$ given
 any multiple of $q_1 q_2$, but we won't show it here. Hint: treat $q_1$ and $q_2$ as unknowns. Construct a system of equations for these unknowns and solve it.


 \begin{solution}
Your solution goes here
\end{solution}
 \end{ppart}


 \begin{ppart} 
Let $k$ be the bit length of $n$. Let $g$ be a generator of $\qrn$. The public key for our hash function is $n$ and $g$. The inputs will be two values in the set $\{1, 2, \dots, q_1 q_2\}$. Let $H(a,b)=g^{a2^k+b}\bmod n$ (note that $a2^k+b$ is
 just concatenation of $a$ and $b$ as bit strings).  Show that this is a
 collision-resistant hash family under the assumption that such $n$ are hard
 to factor. Specifically, show how, given $(a_1, b_1)$ and $(a_2, b_2)$ that
 collide, one can factor $n$.  You may use without proof the fact that $n$
 can be efficiently factored given any non-zero multiple of $q_1q_2$.
 \end{ppart}

 \begin{solution}
Your solution goes here
\end{solution}




 \begin{ppart} 
     Show also that this is a trapdoor hash family (like in Problem 1, you need to find what
 trapdoor information $G$ should output and what the algorithm $T$
 will do).
 \end{ppart}

 \begin{solution}
Your solution goes here
\end{solution}

\end{problem}


\begin{problem} (30 points)
	Show how to compute the root of an $n$-leaf Merkle tree in $O(\log n)$ space.
	More precisely, assume $n$ is a power of 2.  Given data items $x_1, x_2, \dots, x_n$ as values to be placed in the
	leaves of the tree, and a hash function $H$, write pseudocode to compute the
	root $r$ of the tree while never storing more than $2+\log n$
	hash values. Analyze the storage space used by your pseudocode to prove that is never more than $2+\log n$. Note that we do not count the storage of the data items. 
\end{problem}
  
\begin{solution}
Your solution goes here
\end{solution}
\end{document}
